% !TEX root = ../../main.tex

\section{Abstracciones Funcionales}

Las funciones de primera clase, la currificación, la aplicación parcial y las clases de tipos proporcionan la base para las abstracciones estudiadas en esta sección. \texttt{Functor}, \texttt{Applicative} y \texttt{Monad} forman una jerarquía que se distingue por el grado de dependencia que permiten expresar entre las partes de un programa. Un functor permite transformar el resultado de un cómputo; un aplicativo permite combinar cómputos cuya estructura está determinada de antemano; y una mónada permite que un cómputo posterior dependa del valor producido por uno anterior \cite{Wadler1995,McBride2008,Yorgey2009}.

\subsection{Functores}

En \textsf{Haskell}, \texttt{Functor} es una clase de tipos cuyas instancias se definen para constructores de tipos, como las listas \verb|[]| o \verb|Maybe|. Su operación principal es \texttt{fmap}, que permite aplicar una función a los valores asociados a una estructura y obtener un resultado que conserva el mismo constructor de tipos. La definición simplificada de la clase se muestra en el Código~\ref{lst:functor-class}:

\begin{lstlisting}[style=haskell,caption={Definición simplificada de la clase \texttt{Functor}.},label={lst:functor-class}]
class Functor f where
    fmap :: (a -> b) -> f a -> f b
\end{lstlisting}

En esta firma, \verb|f| representa un constructor de tipos, mientras que \verb|a| y \verb|b| representan los tipos de los valores antes y después de aplicar la función. Así, \texttt{fmap} recibe una función de tipo \verb|a -> b| y transforma un valor de tipo \verb|f a| en uno de tipo \verb|f b|, manteniendo el contexto representado por \verb|f|.

En instancias concretas como las listas y \verb|Maybe|, esta operación puede entenderse intuitivamente mediante la metáfora de un contenedor: \texttt{fmap} transforma los valores contenidos sin modificar la estructura que los organiza. Bajo esta interpretación pedagógica, puede imaginarse que la operación accede a los valores, aplica una función sobre ellos y vuelve a representarlos dentro del mismo contexto. Sin embargo, esta metáfora no constituye la definición general de un functor.

Además, toda instancia válida de \texttt{Functor} debe respetar las leyes de identidad y composición. Estas leyes garantizan que la aplicación de \texttt{fmap} preserve de manera coherente la estructura asociada al constructor de tipos \cite{Yorgey2009}.

Los ejemplos clásicos de functores incluyen las listas y los árboles, pues \texttt{fmap} transforma sus elementos de manera uniforme mientras mantiene su estructura subyacente.

Como se mencionó anteriormente, otro ejemplo relevante es \texttt{Maybe}. El tipo \texttt{Maybe a} representa valores opcionales: un valor disponible se expresa mediante el constructor de datos \texttt{Just}, mientras que la ausencia de un valor se representa mediante \texttt{Nothing}. Al utilizar \texttt{fmap} sobre este tipo, la función recibida se aplica al valor contenido en \texttt{Just}, mientras que \texttt{Nothing} permanece sin cambios.

En la implementación de \texttt{base} incluida con GHC, este comportamiento se define mediante la instancia de \texttt{Functor} para el constructor de tipos \texttt{Maybe} mostrada en el Código~\ref{lst:functor-maybe-instance} \cite{GHCInternalBase}:

\begin{lstlisting}[style=haskell,caption={Instancia de \texttt{Functor} para \texttt{Maybe}.},label={lst:functor-maybe-instance}]
instance Functor Maybe where
    fmap _ Nothing  = Nothing
    fmap f (Just a) = Just (f a)
\end{lstlisting}

La primera ecuación establece que, cuando el valor recibido es \texttt{Nothing}, no existe ningún elemento sobre el cual aplicar la función y, por tanto, el resultado permanece como \texttt{Nothing}. El guion bajo indica que la función recibida no necesita utilizarse en este caso. En cambio, cuando el valor tiene la forma \texttt{Just a}, la segunda ecuación aplica la función \texttt{f} al elemento \texttt{a} y vuelve a encapsular el resultado mediante el constructor \texttt{Just}. Ambos comportamientos se ilustran en el Código~\ref{lst:functor-maybe-fmap-examples}:

\begin{lstlisting}[style=haskell,caption={Aplicación de \texttt{fmap} sobre valores de tipo \texttt{Maybe}.},label={lst:functor-maybe-fmap-examples}]
    > fmap (+1) (Just 2)
    Just 3
    > fmap (+1) Nothing
    Nothing
\end{lstlisting}

\subsection{Aplicativos}

En \textsf{Haskell}, \texttt{Applicative} es una clase de tipos definida para constructores de tipos que ya poseen una instancia de \texttt{Functor}. Sus operaciones principales son \texttt{pure} y \verb|(<*>)|: la primera introduce un valor en un contexto, mientras que la segunda aplica una función contextualizada a un valor dentro de ese mismo contexto. De esta forma, la clase caracteriza un estilo aplicativo de programación, menos expresivo que el monádico y, por tanto, aplicable a una mayor variedad de estructuras \cite{McBride2008}. Para ilustrar este comportamiento, se utilizará nuevamente el constructor de tipos \texttt{Maybe} y el operador binario \texttt{(+)}.

Un constructor de tipos que posee una instancia de \texttt{Applicative} se denomina aplicativo. Considerando \texttt{Maybe}, para sumar \texttt{Just 2} y \texttt{Just 3}, la función \texttt{(+)} se introduce primero en el contexto mediante \texttt{pure}. La primera aplicación de \verb|(<*>)| produce la función unaria contextualizada \texttt{Just ((+) 2)}, de tipo \verb|Maybe (Int -> Int)|. La segunda aplicación combina esta función con \texttt{Just 3} y produce el resultado \texttt{Just 5}.

Cada instancia de \texttt{Applicative} define cómo se comportan estas operaciones para un constructor de tipos concreto. En la instancia de \texttt{Maybe}, si alguno de los cómputos asociados a los sumandos produce \texttt{Nothing}, esta ausencia se propaga al resultado. De esta forma, el comportamiento de los valores opcionales queda definido una sola vez en la instancia, evitando repetir ajustes de patrones cada vez que se combinan operaciones que pueden fallar.

Formalmente, la clase \texttt{Applicative} extiende \texttt{Functor} y declara las operaciones \texttt{pure} y \verb|(<*>)| mediante la definición simplificada del Código~\ref{lst:applicative-class}:

\begin{lstlisting}[style=haskell,caption={Definición simplificada de la clase \texttt{Applicative}.},label={lst:applicative-class}]
    class Functor f => Applicative f where
        pure :: a -> f a
        <*>  :: f (a-> b) -> f a -> f b
\end{lstlisting}

En particular, la instancia \verb|Applicative Maybe| define los operadores \verb|pure| y \verb|<*>| de acuerdo con el comportamiento de los valores opcionales: permite introducir valores en el contexto y combinar operaciones que pueden fallar. Una definición de esta instancia se presenta en el Código~\ref{lst:applicative-maybe-instance}:

\begin{lstlisting}[style=haskell,caption={Instancia de \texttt{Applicative} para \texttt{Maybe}.},label={lst:applicative-maybe-instance}]
    instance Applicative Maybe where
        pure = Just
        Nothing <*> _ = Nothing
        _ <*> Nothing = Nothing
        (Just f) <*> (Just x) = Just (f x)
\end{lstlisting}

Nótese que, en caso de operar con algún fallo de cómputo o ausencia de algún valor, cuya abstracción equivale al constructor \verb|Nothing|, las siguientes aplicaciones retornarán a su vez un elemento \verb|Nothing|, propagando así el comportamiento definido por la instancia \verb|Applicative Maybe|. En caso contrario, si ambos operandos se encuentran definidos, es decir, son elementos encapsulados en el constructor \verb|Just|, se procede con la aplicación del primero sobre el segundo, para finalmente empaquetar el elemento resultante dentro del mismo constructor \verb|Just|.

Volviendo al ejemplo de aplicación de la suma \verb|(+)| sobre elementos de tipo \verb|Maybe Int|, se requiere que esta se encuentre encapsulada dentro del constructor \verb|Just|. El operador \verb|<*>| espera que su operando izquierdo sea una función dentro del contexto. Por ello, \verb|pure (+)| introduce la función de suma en \verb|Maybe|, produciendo el valor contextualizado requerido. La expresión aplicativa concreta se muestra en el Código~\ref{lst:applicative-maybe-sum-pure}:

\begin{lstlisting}[style=haskell,caption={Suma aplicativa de dos valores de tipo \texttt{Maybe Int}.},label={lst:applicative-maybe-sum-pure}]
    pure (+) <*> Just 2 <*> Just 3
\end{lstlisting}

Como \texttt{Functor} es una superclase de \texttt{Applicative}, todo constructor de tipos con una instancia de \texttt{Applicative} también dispone de \texttt{fmap}. El operador \verb|(<$>)| corresponde a su forma infija, por lo que se cumple la igualdad \verb|f <$> m = fmap f m|. Además, la coherencia entre ambas clases establece que \verb|fmap f m = pure f <*> m|. En consecuencia, ambas expresiones pueden relacionarse mediante \verb|f <$> m = fmap f m = pure f <*> m|. De este modo, el ejemplo anterior puede escribirse con la sintaxis más directa del Código~\ref{lst:applicative-maybe-sum-infix}:

\begin{lstlisting}[style=haskell,caption={Forma infija de la suma aplicativa sobre \texttt{Maybe}.},label={lst:applicative-maybe-sum-infix}]
    (+) <$> Just 2 <*> Just 3
\end{lstlisting}

Finalmente, nótese que en una expresión aplicativa, la estructura del cómputo queda fijada de antemano: cada argumento se evalúa en su propio contexto y luego se combinan sus resultados mediante la aplicación de una función pura. Esta característica hace explícita la ausencia de dependencias de datos entre los subcómputos asociados a los argumentos, pues ninguno puede usar el resultado del otro para determinar su propia forma. Esto no implica que sus efectos sean conmutativos, aunque puede habilitar una evaluación más paralelizable cuando la implementación lo permite \cite{McBride2008}.

Esta misma restricción implica que, si se quisiera que el segundo argumento dependa del valor producido por el primero, el estilo aplicativo resulta insuficiente, ya que la forma del cómputo no puede adaptarse dinámicamente a resultados intermedios. Por otro lado, la abstracción monádica proporciona el operador \textit{bind} \verb|(>>=)|, que permite construir el segundo cómputo a partir del resultado del primero. Esta capacidad de expresar dependencias entre resultados intermedios motiva la introducción de la clase \texttt{Monad} en \textsf{Haskell} \cite{Wadler1995}.

\subsection{Mónadas}

En \textsf{Haskell}, \texttt{Monad} es una clase de tipos cuyas instancias definen cómo secuenciar cómputos dentro de un mismo contexto cuando un cómputo posterior puede depender del valor producido por uno anterior \cite{Wadler1995}. Como \texttt{Applicative} es una superclase de \texttt{Monad}, todo constructor de tipos que posea una instancia de \texttt{Monad} también debe poseer una instancia de \texttt{Applicative}. El Código~\ref{lst:monad-class} muestra una definición simplificada con las operaciones relevantes:

\begin{lstlisting}[style=haskell,caption={Definición simplificada de la clase \texttt{Monad}.},label={lst:monad-class}]
class Applicative m => Monad m where
    (>>=)  :: m a -> (a -> m b) -> m b
    return :: a -> m a
\end{lstlisting}

Para interpretar la firma de \verb|(>>=)|, consideraremos un cómputo \verb|ma :: m a| y una función \verb|k :: a -> m b|. La expresión \verb|ma >>= k| compone el primer cómputo con la función k de acuerdo con el comportamiento definido por la instancia, por ende entrega su resultado de tipo \verb|a| a la función \verb|k| y obtiene un nuevo cómputo de tipo \verb|m b|. De esta manera, la función \verb|k| puede utilizar el resultado anterior para determinar la forma del cómputo siguiente.

La operación \verb|return| introduce un valor en el contexto monádico y, en la jerarquía actual de clases de tipos, coincide con \verb|pure|. Su nombre no representa una transferencia de control como en lenguajes imperativos, sino la construcción de un valor de tipo \verb|m a| a partir de uno de tipo \verb|a|.

Al encadenar varias aplicaciones de \verb|(>>=)|, las funciones entregadas al operador pueden escribirse como funciones anónimas $\lambda$. Estas permiten asignar nombres a los resultados intermedios y utilizarlos en los cómputos posteriores. Para ilustrar este comportamiento, se retomará el ejemplo de la suma con \verb|Maybe|.

En este caso, la instancia \verb|Monad Maybe| define el comportamiento del operador \verb|>>=| para los valores opcionales. Las operaciones relevantes se presentan en el Código~\ref{lst:monad-maybe-instance}:

\begin{lstlisting}[style=haskell,caption={Instancia de \texttt{Monad} para \texttt{Maybe}.},label={lst:monad-maybe-instance}]
    instance Monad Maybe where
        return x = Just x
        Nothing >>= _ = Nothing
        (Just x) >>= f = f x
\end{lstlisting}

Cuando el primer cómputo produce \verb|Nothing|, la ausencia de valor se propaga como resultado. En cambio, cuando produce un valor de la forma \verb|Just x|, el elemento \verb|x| se entrega como argumento a la función ubicada a la derecha de \verb|>>=|. En este caso, la función que se ejecuta debe respetar la firma de tipos y retornar un nuevo cómputo de tipo \verb|m b|. El Código~\ref{lst:monad-maybe-sum} presenta la aplicación de la función \verb|(+)| mediante notación monádica:

\begin{lstlisting}[style=haskell,caption={Suma monádica de dos valores de tipo \texttt{Maybe}.},label={lst:monad-maybe-sum}]
    Just 3 >>= (\x -> Just 2 >>= (\y -> return ((+) x y)))
\end{lstlisting}

Por tanto, al ejecutar el operador \verb|>>=|, se desempaqueta el número 3 y se aplica la función anónima de la derecha sobre este valor, reemplazando todas las apariciones del identificador \texttt{x} por 3. A continuación, se realiza el mismo proceso con el cómputo \verb|Just 2|, resultando finalmente en la expresión \texttt{((+) 3 2)} que será empaquetada gracias a la instrucción \texttt{return} retornando el valor contextualizado \texttt{Just 5}.

Cabe recalcar que ahora el identificador \verb|x| queda dentro del alcance de la función que construye el segundo cómputo. Por ello, la expresión asociada a \verb|Just 2| puede reemplazarse por otra que utilice \verb|x|. El Código~\ref{lst:monad-maybe-dependent-computation} muestra un programa que utiliza esta nueva expresividad de las mónadas:

\begin{lstlisting}[style=haskell,caption={Cómputo monádico con dependencia entre resultados intermedios.},label={lst:monad-maybe-dependent-computation}]
    Just 3 >>= (\x -> Just ((*) 2 x) >>= (\y -> return ((+) x y)))
\end{lstlisting}

Del programa anterior, el primer cómputo desempaqueta \verb|Just 3| y reemplaza este elemento en las apariciones del identificador \verb|x|. Luego, el cómputo posterior \verb|Just ((*) 2 x)| es capaz de considerar el valor numérico de \verb|x| retornando \verb|Just 6|, luego se desempaqueta su resultado y este reemplaza las apariciones del identificador \verb|y|. Quedando como cómputo final la operación \verb|((+) 3 6)| cuyo resultado termina siendo empaquetado por la instrucción \verb|return| como \verb|Just 9|. Este flujo de ejecución permite una mayor expresividad enriqueciendo el programa. De esta manera la notación monádica termina siendo algo característico y central en \textsf{Haskell}.

Aunque esta composición permita expresar dependencias de forma correcta y versatil, la acumulación de funciones anónimas dificulta su lectura. La \textit{do-notation}, presentada a continuación, proporciona una sintaxis más clara para escribir la misma secuencia monádica.
